---
\documentclass[12pt]{article}

\usepackage{amsmath}
\usepackage{amssymb}
\usepackage{booktabs}
\usepackage{hyperref}
\usepackage{geometry}
\usepackage{lmodern}
\usepackage{microtype}
\usepackage{setspace}
\usepackage{parskip}
\usepackage{xcolor}
\usepackage{listings}

\geometry{margin=1in}
\setstretch{1.15}

\hypersetup{
  colorlinks=true,
  linkcolor=blue!60!black,
  urlcolor=blue!60!black,
  citecolor=blue!60!black
}

\title{TriageCore Resource Estimation Methodology}
\author{Corey}
\date{\today}

\begin{document}

\maketitle
\tableofcontents
\newpage

% ─────────────────────────────────────────────────────────────────────────────
\section{Overview}
% ─────────────────────────────────────────────────────────────────────────────

TriageCore records the estimated resource consumption of every agentic coding
task it orchestrates. Because direct hardware measurement (e.g., RAPL, NVML,
smart-plug readings) is not universally available across desktop environments,
TriageCore uses a \textbf{heuristic estimation model} grounded in the
literature on computational energy accounting. All parameters are configurable
in \texttt{triagecore.toml}; defaults are documented below alongside the
formulae.

Every estimate is stored in the append-only task ledger together with
\texttt{estimation\_method} and \texttt{measurement\_boundary} fields, so that
readers of any derived data can assess the scope and assumptions of each figure.

% ─────────────────────────────────────────────────────────────────────────────
\section{Operational Energy}
% ─────────────────────────────────────────────────────────────────────────────

\subsection*{Formula}

\begin{equation}
  E = \frac{P \cdot t}{3{,}600{,}000}
  \label{eq:energy}
\end{equation}

\begin{table}[h!]
\centering
\begin{tabular}{@{}lll@{}}
\toprule
Symbol & Quantity & Default \\
\midrule
$E$   & Energy consumption (kWh)                    & --- \\
$P$   & Device power draw during inference (W)      & \textbf{300\,W} \\
$t$   & Task wall-clock duration (s)                & measured \\
\bottomrule
\end{tabular}
\caption{Parameters for Equation~\ref{eq:energy}.}
\end{table}

\subsection*{Rationale}

A value of 300\,W is a conservative mid-range estimate for a consumer desktop
GPU (e.g., NVIDIA RTX 3070--4090 class) running local inference at moderate
load. Users with known hardware specifications should override this value in
\texttt{triagecore.toml}. The divisor $3{,}600{,}000$ converts watt-seconds to
kilowatt-hours ($1\,\text{kWh} = 3.6 \times 10^{6}\,\text{J}$).

\textit{Implementation:} \texttt{SustainabilityEstimator.estimate()},
\texttt{triage\_core/sustainability.py}, line~25.

% ─────────────────────────────────────────────────────────────────────────────
\section{Operational Carbon Emissions}
% ─────────────────────────────────────────────────────────────────────────────

\subsection*{Formula}

\begin{equation}
  C_{\text{op}} = E \times I_{\text{grid}}
  \label{eq:emissions}
\end{equation}

\begin{table}[h!]
\centering
\begin{tabular}{@{}lll@{}}
\toprule
Symbol & Quantity & Default \\
\midrule
$C_{\text{op}}$    & Operational carbon emissions (gCO$_2$e)    & --- \\
$E$                & Energy (kWh), from Eq.~\ref{eq:energy}     & --- \\
$I_{\text{grid}}$  & Grid carbon intensity (gCO$_2$e\,kWh$^{-1}$) & \textbf{400} \\
\bottomrule
\end{tabular}
\caption{Parameters for Equation~\ref{eq:emissions}.}
\end{table}

\subsection*{Rationale}

The global average grid carbon intensity is approximately
475\,gCO$_2$e\,kWh$^{-1}$ (IEA, 2023). A default of 400\,gCO$_2$e\,kWh$^{-1}$
represents a modestly cleaner mixed grid (US average
$\approx$\,386\,gCO$_2$e\,kWh$^{-1}$; EU average
$\approx$\,255\,gCO$_2$e\,kWh$^{-1}$). Time-of-day variation is not currently
modelled; this is noted as a limitation (Section~\ref{sec:caveats}).

\textit{Implementation:} line~26.

% ─────────────────────────────────────────────────────────────────────────────
\section{Indirect Water Footprint}
% ─────────────────────────────────────────────────────────────────────────────

\subsection*{Formula}

\begin{equation}
  W = E \times f_{\text{water}}
  \label{eq:water}
\end{equation}

\begin{table}[h!]
\centering
\begin{tabular}{@{}lll@{}}
\toprule
Symbol & Quantity & Default \\
\midrule
$W$                 & Indirect water consumption (litres)       & --- \\
$E$                 & Energy (kWh), from Eq.~\ref{eq:energy}    & --- \\
$f_{\text{water}}$  & Water intensity of electricity (L\,kWh$^{-1}$) & \textbf{1.5} \\
\bottomrule
\end{tabular}
\caption{Parameters for Equation~\ref{eq:water}.}
\end{table}

\subsection*{Rationale}

Electricity generation consumes water through steam-cycle cooling, evaporation
at hydroelectric reservoirs, and fuel processing. The default factor of
$1.5\,\text{L}\,\text{kWh}^{-1}$ is derived from global average estimates for
a mixed thermoelectric--renewable grid \citep{meldrum2013, grubert2018}.
Regional water intensity can vary by more than an order of magnitude.

\textbf{Measurement boundary:} indirect electricity-related water only. Direct
on-premise cooling water is not included.

\textit{Implementation:} line~27.

% ─────────────────────────────────────────────────────────────────────────────
\section{Embodied Carbon Allocation}
% ─────────────────────────────────────────────────────────────────────────────

\subsection*{Formula}

\begin{equation}
  C_{\text{emb}} = C_{\text{device}} \times
    \frac{t\,/\,3600}{H_{\text{lifetime}}}
  \label{eq:embodied}
\end{equation}

\begin{table}[h!]
\centering
\begin{tabular}{@{}lll@{}}
\toprule
Symbol & Quantity & Default \\
\midrule
$C_{\text{emb}}$      & Embodied carbon allocated to task (gCO$_2$e)   & --- \\
$C_{\text{device}}$   & Total device embodied carbon (gCO$_2$e)         & \textbf{300,000} \\
$t$                   & Task duration (s)                               & measured \\
$H_{\text{lifetime}}$ & Expected device lifetime (hours)                & \textbf{20,000} \\
\bottomrule
\end{tabular}
\caption{Parameters for Equation~\ref{eq:embodied}.}
\end{table}

\subsection*{Rationale}

Hardware manufacture accounts for a substantial share of a computing device's
lifetime carbon footprint \citep{gupta2022}. The default of
$300{,}000\,\text{gCO}_2\text{e}$ (300\,kg\,CO$_2$e) is a representative
figure for a mid-range desktop workstation including GPU
\citep{malmodin2018}. A lifetime of 20,000 hours corresponds to approximately
ten years of average daily use ($\sim$5.5\,h/day).

Allocation follows the \textbf{proportional-use method}: the fraction of device
lifetime consumed by a task determines the fraction of embodied carbon
attributed to it. This method is consistent with lifecycle assessment practice
(ISO~14044) and cloud carbon accounting literature \citep{lannelongue2021}.

\textbf{Measurement boundary:} embodied carbon of host device only. Embodied
carbon of model weights (i.e., training compute) is out of scope.

\textit{Implementation:} lines~29--30.

% ─────────────────────────────────────────────────────────────────────────────
\section{Per-Accepted-Task Resource Intensity}
% ─────────────────────────────────────────────────────────────────────────────

\subsection*{Formula}

\begin{equation}
  R_{\text{accepted}} = \frac{\displaystyle\sum_{i} R_i}{N_{\text{accepted}}}
  \label{eq:per_accepted}
\end{equation}

\begin{table}[h!]
\centering
\begin{tabular}{@{}ll@{}}
\toprule
Symbol & Quantity \\
\midrule
$R_{\text{accepted}}$ & Mean resource use per accepted task \\
$\sum R_i$            & Total resource use across all tasks \\
$N_{\text{accepted}}$ & Count of tasks with \texttt{accepted = true} \\
\bottomrule
\end{tabular}
\caption{Parameters for Equation~\ref{eq:per_accepted}.}
\end{table}

This metric is the \textbf{primary scientific outcome variable} of TriageCore
experiments. It penalises configurations that achieve low per-task costs but
high rejection or retry rates, because rejected work still consumed real
resources without producing accepted software artefacts.

\noindent Reported dimensions:
\begin{itemize}
  \item kWh per accepted task
  \item gCO$_2$e per accepted task
  \item Litres of water per accepted task
  \item gCO$_2$e embodied carbon per accepted task
  \item Tokens (input + output) per accepted task
  \item Human review minutes per accepted task
  \item Optional subjective review workload
\end{itemize}

Human review minutes measure elapsed review time. The optional
\texttt{review\_workload} label records the reviewer's perceived assessment
burden (\texttt{not\_recorded}, \texttt{low}, \texttt{medium}, or
\texttt{high}) and should be analysed as a subjective ordinal indicator, not as
an objective performance metric.

\subsection{Local-First Benefit Signals}

The desktop telemetry dashboard intentionally foregrounds local-first benefit
signals because the workbench is also an operator environment. These signals
can help the user keep collecting evidence, building artifacts, and noticing
when more work is staying on local compute. Current dashboard signals include:
\begin{itemize}
  \item accepted yield
  \item percent of tasks routed through local-first runners
  \item count of accepted local-first tasks
  \item percent of tasks that did not require mandatory human review
\end{itemize}

These are motivational and operational indicators, not standalone proof of
environmental savings. Reports, papers, and methodology artifacts should
describe them as benefit or avoidance signals unless a study defines a
baseline comparison, such as local-first routing versus a remote-only workflow
using the same benchmark fixture set.

% ─────────────────────────────────────────────────────────────────────────────
\section{Model and Backend Comparison}
% ─────────────────────────────────────────────────────────────────────────────

Model/backend comparison studies use the same benchmark fixture set across
each candidate configuration and tag all evidence with a shared
\texttt{study\_id} plus a unique \texttt{run\_id} for each backend/model pair.
Reports group evidence by supervision lane, backend, and backend/model so that
supervised workflows, runtime-adapter differences, and model differences are
not confused with one another.

Comparison runs should use:
\begin{itemize}
  \item identical benchmark fixtures
  \item identical timeout settings unless timeout is the experimental variable
  \item explicit backend type, base URL when applicable, and model identifier
  \item separate \texttt{run\_id} values for every backend/model pair
  \item the combined \texttt{benchmark-report --study-id <study>} output for
    interpretation
  \item the \texttt{By Supervision} section when local-only and supervised
    outcomes are mixed
\end{itemize}

Expected destructive-task handoff is treated as correct safety behavior. A
configuration should only be preferred when it improves accepted outcomes
without increasing unexpected handoffs, validator failures, or subjective review
burden.

% ─────────────────────────────────────────────────────────────────────────────
\section{Supervised Hybrid Workflows}
% ─────────────────────────────────────────────────────────────────────────────

TriageCore distinguishes local-only execution from supervised hybrid workflows.
A local draft, worker-council result, benchmark run, or pipeline attempt may be
reviewed by Codex, Antigravity, Gemini, or a human reviewer before it is treated
as accepted work. These reviews are recorded as
\texttt{supervisor\_reviewed} events rather than merged into the original local
execution event.

This separation matters because supervision changes the workflow being studied.
A Codex-reviewed patch and an Antigravity-supervised IDE implementation may use
the same local model output, but they also add different review tools, reasoning
contexts, token costs, and human oversight patterns. Reports should therefore
label supervised outcomes separately from local-only outcomes whenever they are
used as scientific evidence.

Supervisor review fields include the supervising tool, model, profile,
decision, notes, linked artifact path, and estimated supervisor input/output
tokens when exact usage is not available. Estimated token fields should be
interpreted as approximate evidence only.

Benchmark reports summarize supervisor reviews under the same
\texttt{study\_id} and \texttt{run\_id} filters as the benchmark evidence. The
\texttt{Supervisor Reviews} table reports review counts, decision counts, and
estimated supervisor token totals by tool. Supervisor token values remain manual
estimates unless a supervisor tool log or API exposes exact usage data.

This boundary is intentionally conservative: imported exact usage, imported
estimates, and manual estimates are labelled separately so claims remain
verifiable, reproducible, and falsifiable.

% ─────────────────────────────────────────────────────────────────────────────
\section{Default Parameter Summary}
% ─────────────────────────────────────────────────────────────────────────────

\begin{table}[h!]
\centering
\begin{tabular}{@{}lllll@{}}
\toprule
Parameter & Symbol & Default & Unit & \texttt{triagecore.toml} key \\
\midrule
Device power draw    & $P$                  & 300       & W               & \texttt{sustainability.default\_watts} \\
Grid intensity       & $I_{\text{grid}}$    & 400       & gCO$_2$e\,kWh$^{-1}$ & \texttt{sustainability.grid\_intensity\_gco2e\_per\_kwh} \\
Water intensity      & $f_{\text{water}}$   & 1.5       & L\,kWh$^{-1}$   & \texttt{sustainability.water\_intensity\_l\_per\_kwh} \\
Device embodied C    & $C_{\text{device}}$  & 300,000   & gCO$_2$e        & \texttt{sustainability.device\_embodied\_gco2e} \\
Device lifetime      & $H_{\text{lifetime}}$& 20,000    & hours           & \texttt{sustainability.device\_lifetime\_hours} \\
\bottomrule
\end{tabular}
\caption{Default parameters and their \texttt{triagecore.toml} override keys.}
\end{table}

% ─────────────────────────────────────────────────────────────────────────────
\section{Estimation Caveats and Limitations}
\label{sec:caveats}
% ─────────────────────────────────────────────────────────────────────────────

\begin{enumerate}
  \item \textbf{Power draw is assumed, not measured.} Direct measurement via
    RAPL (Intel/AMD), NVML (NVIDIA), or a smart plug would substantially
    improve accuracy. The current model assumes constant load at the configured
    wattage for the full task duration; idle periods within a task are not
    accounted for.

  \item \textbf{Grid intensity is static.} Carbon intensity varies by time of
    day and season. Tools such as ElectricityMaps or the WattTime API could
    provide real-time intensity values as a future enhancement.

  \item \textbf{Water estimates are regional-average proxies.} The factor of
    $1.5\,\text{L}\,\text{kWh}^{-1}$ may underestimate intensity for coal-heavy
    grids or overestimate for predominantly hydro/wind grids.

  \item \textbf{Embodied carbon defaults are order-of-magnitude estimates.}
    Actual figures depend on hardware manufacturer, manufacturing region, and
    device class. Boavizta device profiles or manufacturer-published product
    carbon footprint disclosures should be used where available.

  \item \textbf{Training carbon is out of scope.} The embodied carbon of the
    model weights (i.e., the compute used to train the LLM) is not attributed
    to inference tasks in this version. This is a deliberate boundary choice,
    consistent with operational-boundary carbon accounting, and should be stated
    explicitly in any published results.

  \item \textbf{All estimates are labelled as estimates.} Every ledger record
    includes
    \texttt{"estimation\_method": "heuristic\_profile\_v1"} and
    \texttt{"measurement\_boundary": "local operational estimate plus optional
    embodied amortisation"} to signal this scope to downstream consumers.
\end{enumerate}

% ─────────────────────────────────────────────────────────────────────────────
\section{Execution Hardware Specification}
% ─────────────────────────────────────────────────────────────────────────────

For the comparison study of local model execution and specialist councils (Study 002), the benchmarking environment was standardized on a host workstation with the following hardware specifications:

\begin{itemize}
  \item \textbf{Host System}: Alienware m18 R1 AMD Laptop
  \item \textbf{Processor (CPU)}: AMD Ryzen 9 7845HX with Radeon Graphics (12 physical cores, 24 threads, AVX/AVX2 support)
  \item \textbf{System Memory (RAM)}: 32\,GB DDR5 RAM (33,487,876,096 bytes)
  \item \textbf{Dedicated Graphics (GPU)}: NVIDIA GeForce RTX 4070 Laptop GPU (Discrete, Compute Capability 8.9, 8\,GB VRAM / 8,585,216,000 bytes)
  \item \textbf{Operating System}: Windows 11 Pro 64-bit (Build 26200)
\end{itemize}

This system configuration is representative of developer-class workstations where local model offloading is actively deployed to protect the cloud compute budget.

% ─────────────────────────────────────────────────────────────────────────────
\section*{References}
% ─────────────────────────────────────────────────────────────────────────────

\begin{thebibliography}{9}

\bibitem{grubert2018}
Grubert, E., \& Sanders, K.~T. (2018).
Water use in the United States energy system: A national assessment and unit
process inventory of water consumption and withdrawals.
\textit{Environmental Science \& Technology, 52}(11), 6695--6703.
\url{https://doi.org/10.1021/acs.est.8b00139}

\bibitem{gupta2022}
Gupta, U., Kim, Y.~G., Lee, S., Tse, J., Lee, H.~H.~S., Wei, G.-Y.,
Brooks, D., \& Wu, C.-J. (2022).
Chasing carbon: The elusive environmental footprint of computing.
\textit{IEEE Micro, 42}(4), 37--47.
\url{https://doi.org/10.1109/MM.2022.3163226}

\bibitem{iea2023}
International Energy Agency. (2023).
\textit{Electricity 2024: Analysis and forecast to 2026}.
IEA. \url{https://www.iea.org/reports/electricity-2024}

\bibitem{lannelongue2021}
Lannelongue, L., Grealey, J., \& Inouye, M. (2021).
Green algorithms: Quantifying the carbon footprint of computation.
\textit{Advanced Science, 8}(12), 2100707.
\url{https://doi.org/10.1002/advs.202100707}

\bibitem{malmodin2018}
Malmodin, J., \& Lund\'{e}n, D. (2018).
The energy and carbon footprint of the global ICT and E\&M sectors 2010--2015.
\textit{Sustainability, 10}(9), 3027.
\url{https://doi.org/10.3390/su10093027}

\bibitem{meldrum2013}
Meldrum, J., Nettles-Anderson, S., Heath, G., \& Macknick, J. (2013).
Life cycle water use for electricity generation: A review and harmonization of
literature estimates.
\textit{Environmental Research Letters, 8}(1), 015031.
\url{https://doi.org/10.1088/1748-9326/8/1/015031}

\end{thebibliography}

\end{document}
